\section{Vorwort}
\label{sec:vorwort}
Mit dem Thema Barrierefreiheit bin ich zum ersten Mal während meines Studiums im Modul Webentwicklung in Berührung gekommen. Seitdem hat dieses Thema einen besonderen Stellenwert für mich.

Bei unseren Anwendungen mussten wir immer besonders auf die semantische Struktur achten. Damit war ein Bereich der Barrierefreiheit abgedeckt. Durch viele Negativbeispiele wurde mir die Bedeutung der Barrierefreiheit bewusst. Durch fehlende oder eingeschränkte Barrierefreiheit schließt man einen Teil der Gesellschaft aus.

Das Thema hat mich von Anfang an interessiert, da es mir Spaß gemacht hat, Möglichkeiten zu finden, Anwendungen barrierefreier zu gestalten. In den jeweiligen Modulen gab es dafür auch viel Unterstützung und wurde teilweise sogar vorausgesetzt. 

Als ich mich dann neben dem Studium mit der Entwicklung von Webanwendungen selbstständig machte, blieb das Thema ein weiterer Bestandteil meiner Arbeit. Die Projekte, die ich für meine Kunden entwickelte, sollten weiterhin möglichst vielen Menschen zugänglich gemacht werden. Teilweise wurde dies sogar von den Kunden gefordert, was die Sache natürlich noch wichtiger machte. Schon damals war ich stolz darauf, dass meine Anwendungen barrierefrei waren und diverse Tests auf Barrierefreiheit bestanden hatten.

\section{Motivation} 
\label{sec:motivation}
Dabei habe ich aber nur meine Grundkenntnisse angewandt und gehofft, dass das ausreicht. Im Großen und Ganzen hat es für den Anfang auch gereicht. Es gab aber nie genaue Vorgaben, eher nur Richtlinien, an die man sich aber nicht halten musste. 

Barrierefreiheit war ein viel belächeltes Thema, das aber nur von wenigen ernsthaft umgesetzt wurde. Das Thema ist bei Entwicklern sehr unbeliebt, weil es mit viel Aufwand und Bürokratie verbunden ist. Während meines Studiums und meiner freiberuflichen Tätigkeit habe ich jedoch gemerkt, dass man Barrierefreiheit sehr gut in seine Arbeit integrieren kann. Durch die fehlenden Vorgaben haben sich viele Entwickler aus der Verantwortung gestohlen, ihre Anwendungen barrierefrei zu gestalten.

Dank der EU wurde eine Richtlinie auf den Weg gebracht, die in Deutschland durch das BFSG vorgeschrieben ist. Doch was bedeuten diese Vorgaben nun für die Entwickler, die sie in ihren Anwendungen umsetzen müssen? Wie bereits erwähnt, ist Barrierefreiheit in Entwicklerkreisen kein populäres Thema. Daher war es mein Ziel, mit dieser Arbeit eine Brücke zwischen Entwicklern und diesem Gesetz zu schlagen. Und natürlich meinen Beitrag dazu zu leisten, Anwendungen für möglichst viele Menschen zugänglich zu machen. 

Was bedeutet das Gesetz für mich? Was muss ich beachten, wenn ich eine neue Anwendung entwickle? Sind diese Anforderungen mit modernen Webtechnologien überhaupt umsetzbar? Diese und andere Fragen versuche ich in meiner Arbeit zu beantworten.

\vspace{1em}
\noindent
Zusammengefasst unter der Thesis Frage:
\vspace{1em}

\noindent\textbf{\large Wie lassen sich die Anforderungen des Barrierefreiheitsstärkungsgesetzes (BFSG) in einer modernen Webanwendung technisch umsetzen?}
\vspace{1em}